\documentclass[12pt,a4paper]{article}

% ============================================================
% PACKAGES
% ============================================================
\usepackage[utf8]{inputenc}
\usepackage[T1]{fontenc}
\usepackage[margin=2.5cm]{geometry}
\usepackage{array}
\usepackage{booktabs}
\usepackage{longtable}
\usepackage[table,dvipsnames]{xcolor}
\usepackage[hidelinks,colorlinks=true]{hyperref}
\usepackage{enumitem}
\usepackage{titlesec}
\usepackage{fancyhdr}
\usepackage[most]{tcolorbox}
\usepackage{tabularx}
\usepackage{setspace}
\usepackage{parskip}
\usepackage{graphicx}

% ============================================================
% COLOR PALETTE
% ============================================================
\definecolor{accentdark}{HTML}{0D4F6B}
\definecolor{accentlight}{HTML}{E6F2F8}
\definecolor{principlebg}{HTML}{FFF6E6}
\definecolor{principleborder}{HTML}{C4900A}
\definecolor{grayrule}{HTML}{CCCCCC}
\definecolor{alertbg}{HTML}{FFF0F0}
\definecolor{alertborder}{HTML}{C0392B}

% ============================================================
% HYPERREF
% ============================================================
\hypersetup{
    linkcolor=accentdark,
    urlcolor=accentdark,
    citecolor=accentdark,
    pdftitle={GenAI-Assisted Document Redaction Governance},
    pdfauthor={Andrea Gennari}
}

% ============================================================
% SECTION STYLING
% ============================================================
\titleformat{\section}
  {\Large\bfseries\color{accentdark}}
  {\thesection}{1em}{}
  [\vspace{2pt}{\color{accentdark}\titlerule[1.2pt]}]
\titlespacing*{\section}{0pt}{24pt plus 4pt minus 2pt}{10pt plus 2pt minus 2pt}

\titleformat{\subsection}
  {\large\bfseries\color{accentdark!85!black}}
  {\thesubsection}{0.8em}{}
\titlespacing*{\subsection}{0pt}{18pt plus 3pt minus 2pt}{6pt plus 2pt minus 1pt}

% ============================================================
% HEADER / FOOTER
% ============================================================
\setlength{\headheight}{16pt}
\addtolength{\topmargin}{-4pt}
\pagestyle{fancy}
\fancyhf{}
\fancyhead[L]{\small\color{accentdark!70!black}\textit{GenAI-Assisted Document Redaction Governance}}
\fancyhead[R]{\small\color{accentdark!70!black}\textit{\nouppercase{\rightmark}}}
\fancyfoot[C]{\small\color{accentdark!70!black}\thepage}
\renewcommand{\headrulewidth}{0.4pt}
\renewcommand{\headrule}{\hbox to\headwidth{\color{accentdark!40!white}\leaders\hrule height \headrulewidth\hfill}}
\renewcommand{\footrulewidth}{0pt}

% ============================================================
% PARAGRAPH & LIST SPACING
% ============================================================
\setlength{\parskip}{6pt plus 1pt minus 1pt}
\setlength{\parindent}{0pt}

\setlist[itemize]{
  topsep=4pt, itemsep=3pt, parsep=1pt,
  leftmargin=1.8em,
  label=\textcolor{accentdark}{\textbullet}
}
\setlist[enumerate]{
  topsep=4pt, itemsep=3pt, parsep=1pt,
  leftmargin=2.2em,
  label=\textcolor{accentdark}{\arabic*.}
}

% ============================================================
% TCOLORBOX STYLES
% ============================================================
\newtcolorbox{definitionbox}[1][]{%
  enhanced, breakable,
  colback=accentlight,
  colframe=accentdark,
  coltitle=white,
  fonttitle=\bfseries,
  boxrule=0.8pt,
  arc=2pt,
  left=8pt, right=8pt, top=6pt, bottom=6pt,
  title={#1},
  attach boxed title to top left={yshift=-2mm, xshift=6mm},
  boxed title style={colback=accentdark, arc=2pt, boxrule=0pt},
}

\newtcolorbox{principlebox}[1][]{%
  enhanced, breakable,
  colback=principlebg,
  colframe=principleborder,
  coltitle=principleborder!80!black,
  fonttitle=\bfseries,
  boxrule=0.8pt,
  arc=2pt,
  left=8pt, right=8pt, top=6pt, bottom=6pt,
  title={#1},
  attach boxed title to top left={yshift=-2mm, xshift=6mm},
  boxed title style={colback=principlebg, arc=2pt, boxrule=0.6pt, colframe=principleborder},
}

\newtcolorbox{alertbox}[1][]{%
  enhanced, breakable,
  colback=alertbg,
  colframe=alertborder,
  coltitle=white,
  fonttitle=\bfseries,
  boxrule=0.8pt,
  arc=2pt,
  left=8pt, right=8pt, top=6pt, bottom=6pt,
  title={#1},
  attach boxed title to top left={yshift=-2mm, xshift=6mm},
  boxed title style={colback=alertborder, arc=2pt, boxrule=0pt},
}

\newtcolorbox{quotebox}{%
  enhanced, breakable,
  colback=accentlight!40!white,
  colframe=accentlight!40!white,
  boxrule=0pt,
  borderline west={2.5pt}{0pt}{accentdark},
  arc=0pt,
  left=10pt, right=8pt, top=6pt, bottom=6pt,
  fontupper=\itshape,
}

% ============================================================
% DOCUMENT BEGIN
% ============================================================
\begin{document}

% ============================================================
% TITLE PAGE
% ============================================================
\begin{titlepage}
    \centering
    \vspace*{3cm}

    {\color{accentdark}\rule{\textwidth}{1.5pt}}

    \vspace{1.2cm}
    {\Huge\bfseries\color{accentdark} GenAI-Assisted Document\par}
    \vspace{0.3cm}
    {\Huge\bfseries\color{accentdark} Redaction Governance\par}
    \vspace{0.8cm}
    {\Large\color{accentdark!80!black} Privacy, Compliance and Data Leakage Risk Management\par}
    \vspace{1.2cm}

    {\color{accentdark}\rule{0.4\textwidth}{0.8pt}}

    \vspace{0.8cm}
    \begin{tcolorbox}[
      enhanced,
      colback=alertbg,
      colframe=alertborder,
      arc=3pt,
      boxrule=0.8pt,
      left=12pt, right=12pt, top=8pt, bottom=8pt,
      width=0.85\textwidth,
    ]
      \centering
      {\large\bfseries\color{alertborder} Case Study:\par}
      \vspace{4pt}
      {\large\color{alertborder!80!black} Post Office Horizon Data Breach (2024)\par}
    \end{tcolorbox}

    \vspace{2cm}
    {\large\color{accentdark!70!black} Submitted by\par}
    \vspace{0.3cm}
    {\LARGE\bfseries\color{accentdark} Andrea Gennari\par}
    \vspace{0.5cm}
    {\normalsize\color{accentdark!70!black} Business and Project Management\par}
    \vfill
    {\color{accentdark}\rule{\textwidth}{1.5pt}}
    \vspace{0.5cm}
    {\large\color{accentdark!70!black} \today\par}
\end{titlepage}

% ============================================================
% ABSTRACT
% ============================================================
\thispagestyle{empty}
\vspace*{2cm}
\begin{tcolorbox}[
  enhanced,
  colback=accentlight!30!white,
  colframe=accentdark,
  coltitle=white,
  fonttitle=\bfseries\large,
  title=Abstract,
  attach boxed title to top center={yshift=-3mm},
  boxed title style={colback=accentdark, arc=3pt, boxrule=0pt},
  arc=3pt,
  boxrule=0.8pt,
  left=12pt, right=12pt, top=10pt, bottom=10pt,
]
This project analyses how organizations can safely adopt Generative AI-assisted document redaction tools while managing privacy, compliance and data leakage risks. Using the Post Office Horizon data breach of 2024 as a case study, it proposes a governance framework covering human-in-the-loop controls, audit logging, risk management, RACI responsibilities and measurable KPIs. The central contribution is not a new AI product but a managerial and organizational framework for safe, auditable and compliant adoption of GenAI-assisted redaction in document publication workflows.
\end{tcolorbox}

\newpage
\tableofcontents
\newpage

% ============================================================
% SECTION 1
% ============================================================
\section{Project Overview}

This project analyses how organizations can safely adopt Generative AI-assisted document redaction tools while managing privacy, compliance and data leakage risks.

The project does not aim to create a new commercial redaction tool. Existing commercial solutions already provide AI-assisted or automated redaction capabilities. Instead, the project focuses on the managerial and organizational challenge: how to implement such tools in a controlled, auditable and compliant way.

The proposed case study is the Post Office Horizon data breach. In 2024, Post Office Limited mistakenly uploaded an unredacted version of a legal settlement document to its corporate website. The document exposed personal data of hundreds of people involved in litigation related to the Horizon scandal. The Information Commissioner's Office later found that the organization had failed to implement appropriate technical and organizational measures.

% ============================================================
% SECTION 2
% ============================================================
\section{Research Question}

\begin{quotebox}
How can organizations safely adopt GenAI-assisted document redaction through governance, human-in-the-loop controls, audit logging and risk management?
\end{quotebox}

\medskip

A secondary research question is:

\begin{quotebox}
How could a GenAI-assisted governance framework have reduced the privacy and data leakage risks in the Post Office Horizon document publication case?
\end{quotebox}

% ============================================================
% SECTION 3
% ============================================================
\section{Reason for Choosing the Case}

The Post Office Horizon case is relevant because it shows that document redaction failures are not only technical problems. They are also governance, process and accountability failures.

The case is suitable for a Business and Project Management project because it involves:

\begin{itemize}
    \item privacy and compliance risk;
    \item publication of a wrong document version;
    \item lack of sufficient technical and organizational controls;
    \item reputational damage;
    \item need for clear roles, approvals and auditability;
    \item opportunity to evaluate how GenAI can support, but not replace, human decision-making.
\end{itemize}

% ============================================================
% SECTION 4
% ============================================================
\section{Scope of the Project}

\subsection{In Scope}

\begin{definitionbox}[In Scope]
The project will cover:
\begin{itemize}
    \item analysis of the Post Office Horizon data breach as a case study;
    \item identification of root causes from a managerial perspective;
    \item design of a GenAI-assisted redaction governance framework;
    \item definition of roles and responsibilities;
    \item human-in-the-loop workflow design;
    \item audit log requirements;
    \item risk register;
    \item KPIs;
    \item business case for adoption.
\end{itemize}
\end{definitionbox}

\subsection{Out of Scope}

\begin{tcolorbox}[
  enhanced, breakable,
  colback=gray!6,
  colframe=gray!40,
  boxrule=0.8pt,
  arc=2pt,
  left=8pt, right=8pt, top=6pt, bottom=6pt,
  title={\bfseries Out of Scope},
  coltitle=gray!60!black,
  attach boxed title to top left={yshift=-2mm, xshift=6mm},
  boxed title style={colback=gray!15, arc=2pt, boxrule=0.6pt, colframe=gray!40},
]
The project will not cover:
\begin{itemize}
    \item development of a full commercial software platform;
    \item training of a proprietary AI model;
    \item legal advice on GDPR compliance;
    \item implementation in a real company;
    \item complete technical architecture of an enterprise IT system.
\end{itemize}
\end{tcolorbox}

% ============================================================
% SECTION 5
% ============================================================
\section{Case Background}

\begin{alertbox}[Post Office Horizon Data Breach -- 2024]
Post Office Limited published an unredacted legal settlement document on its corporate website. The document exposed personal data of hundreds of individuals involved in Horizon-related group litigation. The Information Commissioner's Office later found that the organization had failed to implement appropriate technical and organizational measures.
\end{alertbox}

\medskip

The Post Office Horizon scandal involved serious organizational, legal and technological failures. For this project, the case will be used as an example of how a document publication process can fail when there are insufficient controls over redaction, approval, version management and accountability.

The data breach involved:

\begin{itemize}
    \item publication of an unredacted legal settlement document;
    \item exposure of personal data;
    \item affected individuals involved in Horizon-related group litigation;
    \item public availability of the document before removal;
    \item regulatory scrutiny by the Information Commissioner's Office.
\end{itemize}

% ============================================================
% SECTION 6
% ============================================================
\section{Problem Statement}

Organizations often need to share or publish documents with external stakeholders, such as consultants, auditors, suppliers, clients, regulators or the general public. These documents may contain sensitive information, including personal data, addresses, financial information, contractual clauses or commercially sensitive information.

Manual redaction processes are often slow, inconsistent and dependent on individual attention. This creates the risk of:

\begin{itemize}
    \item personal data leakage;
    \item regulatory non-compliance;
    \item reputational damage;
    \item legal exposure;
    \item loss of trust;
    \item inability to prove who approved a document and why.
\end{itemize}

\begin{principlebox}[Central Problem]
The central problem is therefore not only \textit{how to redact documents}, but \textbf{how to govern the redaction process}.
\end{principlebox}

% ============================================================
% SECTION 7
% ============================================================
\section{Proposed Solution}

The proposed solution is a GenAI-assisted document redaction governance framework.

The framework uses GenAI as a support mechanism for identifying sensitive information, assigning risk levels and suggesting redactions. However, the final decision remains under human control.

\begin{principlebox}[Guiding Principle]
\centering
\large\textbf{AI suggests, humans approve, the organization remains accountable.}
\end{principlebox}

% ============================================================
% SECTION 8
% ============================================================
\section{Proposed Workflow}

\begin{definitionbox}[11-Step Governance Workflow]
\begin{enumerate}
    \item Document selected for sharing or publication.
    \item Version control check.
    \item GenAI scan for sensitive data.
    \item Document classification.
    \item Risk score assignment.
    \item AI-generated redaction proposal.
    \item Human review.
    \item Escalation to Legal, DPO or Manager if high-risk.
    \item Final approval.
    \item Audit log generation.
    \item Secure sharing or publication.
\end{enumerate}
\end{definitionbox}

% ============================================================
% SECTION 9
% ============================================================
\section{Human-in-the-Loop Model}

The system should not allow automatic publication of high-risk documents. Human review is mandatory, especially when the document contains personal data, financial information, legal claims, confidential clauses or sensitive internal information.

\begin{table}[h!]
\centering
\rowcolors{2}{accentlight!30!white}{white}
\begin{tabular}{@{} >{\bfseries\color{accentdark}}p{3.5cm} p{3.5cm} p{5.5cm} @{}}
\toprule
\textbf{\color{accentdark}Risk Level} & \textbf{\color{accentdark}Automation Level} & \textbf{\color{accentdark}Required Human Control} \\
\midrule
Low & AI-assisted & Document owner review \\
Medium & AI-assisted with warning & Document owner approval \\
High & AI-assisted only & Legal, DPO or manager approval \\
Very High & No automatic release & Mandatory DPO and senior approval \\
\bottomrule
\end{tabular}
\caption{Human-in-the-loop control by document risk level}
\end{table}

% ============================================================
% SECTION 10
% ============================================================
\section{Roles and Responsibilities}

A RACI matrix defines responsibilities across the governance workflow.

\begin{longtable}{@{} >{\bfseries\color{accentdark}}p{4.2cm} >{\centering\arraybackslash}p{2cm} >{\centering\arraybackslash}p{1.8cm} >{\centering\arraybackslash}p{2cm} >{\centering\arraybackslash}p{2cm} >{\centering\arraybackslash}p{1.8cm} @{}}
\toprule
\textbf{\color{accentdark}Activity} & \textbf{\color{accentdark}Document Owner} & \textbf{\color{accentdark}Legal} & \textbf{\color{accentdark}DPO/ Privacy} & \textbf{\color{accentdark}IT Security} & \textbf{\color{accentdark}Manager} \\
\midrule
\rowcolor{accentlight!30!white}
Document upload & R & C & I & I & A \\
Version control check & R & C & I & C & A \\
\rowcolor{accentlight!30!white}
Sensitive data scan & R & C & C & A & I \\
Redaction proposal review & R & C & C & I & A \\
\rowcolor{accentlight!30!white}
High-risk approval & C & A & A & C & R \\
Final publication & R & I & I & C & A \\
\rowcolor{accentlight!30!white}
Audit review & C & C & A & R & I \\
\bottomrule
\caption{RACI matrix for GenAI-assisted document redaction governance}
\end{longtable}

\begin{tcolorbox}[
  enhanced,
  colback=gray!6,
  colframe=gray!35,
  arc=2pt,
  boxrule=0.6pt,
  left=8pt, right=8pt, top=5pt, bottom=5pt,
]
\textbf{Legend:}\quad
R\,=\,Responsible\quad
A\,=\,Accountable\quad
C\,=\,Consulted\quad
I\,=\,Informed
\end{tcolorbox}

% ============================================================
% SECTION 11
% ============================================================
\section{Risk Register}

\begin{longtable}{@{} p{4cm} >{\centering\arraybackslash}p{2cm} >{\centering\arraybackslash}p{2cm} p{5.5cm} @{}}
\toprule
\textbf{\color{accentdark}Risk} & \textbf{\color{accentdark}Probability} & \textbf{\color{accentdark}Impact} & \textbf{\color{accentdark}Mitigation} \\
\midrule
\rowcolor{accentlight!30!white}
Sensitive data not detected by AI & Medium & High & Mandatory human review for high-risk documents \\
Wrong document version uploaded & Medium & High & Version control and pre-publication checklist \\
\rowcolor{accentlight!30!white}
Overreliance on AI & Medium & High & Policy stating that AI is only a support tool \\
Lack of auditability & Medium & High & Mandatory audit log for each document \\
\rowcolor{accentlight!30!white}
Data exposure to external AI vendor & Medium & High & Vendor assessment and data processing agreement \\
False positives in redaction & High & Medium & Human reviewer can modify AI suggestions \\
\rowcolor{accentlight!30!white}
Lack of staff training & Medium & Medium & Training before implementation \\
Unclear accountability & Medium & High & RACI matrix and approval workflow \\
\bottomrule
\caption{Risk register}
\end{longtable}

% ============================================================
% SECTION 12
% ============================================================
\section{Audit Log Requirements}

An audit log is necessary to prove that the organization followed a controlled process before sharing or publishing a document.

\begin{definitionbox}[Required Audit Log Fields]
\begin{itemize}
    \item document ID;
    \item document type;
    \item upload date and time;
    \item user who uploaded the document;
    \item intended audience;
    \item sensitive data categories detected;
    \item risk score;
    \item AI redaction suggestions;
    \item human reviewer;
    \item changes made by the reviewer;
    \item final approver;
    \item approval date and time;
    \item final publication or sharing status.
\end{itemize}
\end{definitionbox}

% ============================================================
% SECTION 13
% ============================================================
\section{Key Performance Indicators}

\begin{longtable}{@{} p{5cm} p{5.5cm} >{\centering\arraybackslash}p{3cm} @{}}
\toprule
\textbf{\color{accentdark}KPI} & \textbf{\color{accentdark}Formula} & \textbf{\color{accentdark}Target} \\
\midrule
\rowcolor{accentlight!30!white}
Pre-publication scan rate & Documents scanned / documents published & 100\% \\
High-risk approval rate & High-risk docs approved by DPO or Legal / total high-risk docs & 100\% \\
\rowcolor{accentlight!30!white}
Audit completeness & Documents with complete audit log / total processed & 100\% \\
False negative rate & Sensitive data not detected / total sensitive items & $<5\%$ in pilot \\
\rowcolor{accentlight!30!white}
Wrong-version incidents & Number of wrong document versions published & 0 \\
Average review time & Average time from upload to approval & $<24$ hours \\
\rowcolor{accentlight!30!white}
User adoption rate & Active users / target users & $>70\%$ after pilot \\
Data leakage incidents & Number of confirmed leakage incidents & 0 during pilot \\
\bottomrule
\caption{KPIs for the proposed governance framework}
\end{longtable}

% ============================================================
% SECTION 14
% ============================================================
\section{Business Case}

The business case is mainly risk-driven rather than cost-saving driven.

The adoption of a GenAI-assisted redaction governance framework may not immediately reduce costs, especially if software licenses, integration and training are required. However, the main value comes from reducing the probability and impact of:

\begin{itemize}
    \item data breaches;
    \item regulatory action;
    \item reputational damage;
    \item compensation costs;
    \item operational disruption;
    \item loss of trust among stakeholders.
\end{itemize}

\subsection{Example Cost Scenario}

Assume an organization reviews 200 documents per year before external sharing or publication.

\begin{table}[h!]
\centering
\rowcolors{2}{accentlight!30!white}{white}
\begin{tabular}{@{} l r r @{}}
\toprule
\textbf{\color{accentdark}Item} & \textbf{\color{accentdark}Manual Process} & \textbf{\color{accentdark}GenAI-Assisted Process} \\
\midrule
Documents per year & 200 & 200 \\
Average review time & 45 minutes & 20 minutes \\
Total review hours & 150 & 66.7 \\
Hourly cost & EUR 35 & EUR 35 \\
Human review cost & EUR 5,250 & EUR 2,334 \\
Software/API cost & EUR 0 & EUR 10,000 \\
Total annual cost & EUR 5,250 & EUR 12,334 \\
\bottomrule
\end{tabular}
\caption{Illustrative business case}
\end{table}

\begin{principlebox}[Key Insight]
Although the GenAI-assisted process may appear more expensive in direct annual costs, the value lies in risk reduction, auditability and prevention of potentially severe privacy incidents.
\end{principlebox}

% ============================================================
% SECTION 15
% ============================================================
\section{Expected Contribution}

The expected contribution of the project is not a new AI product, but a managerial framework for safe adoption.

The project contributes by:

\begin{itemize}
    \item connecting GenAI adoption with privacy and compliance risk management;
    \item showing how human-in-the-loop controls can reduce automation risk;
    \item defining clear organizational responsibilities;
    \item proposing audit log requirements;
    \item identifying measurable KPIs;
    \item applying the framework to a real-world data breach case.
\end{itemize}

% ============================================================
% PART II — RQ1 IMPLEMENTATION (SKELETON)
% Build order: skeleton first, then fill section by section.
% Each subsection carries a % FILL: note describing intended content.
% Remove the orientation note + [To be developed] stubs before submission.
% ============================================================
\clearpage
\begin{center}
{\color{accentdark}\rule{\textwidth}{1.2pt}}\\[6pt]
{\Huge\bfseries\color{accentdark} Part II}\\[8pt]
{\LARGE\bfseries\color{accentdark!85!black} RQ1 --- A GenAI-Assisted Redaction Governance Framework}\\[8pt]
{\color{accentdark}\rule{\textwidth}{1.2pt}}
\end{center}
\vspace{8pt}

\begin{quotebox}
\textbf{RQ1.} How can organizations safely adopt GenAI-assisted document redaction through governance, human-in-the-loop controls, audit logging and risk management?
\end{quotebox}

% --- Orientation note for the author (remove before final submission) ---
\begin{tcolorbox}[
  enhanced, breakable, colback=gray!6, colframe=gray!40, boxrule=0.6pt, arc=2pt,
  left=8pt, right=8pt, top=5pt, bottom=5pt,
  title={\bfseries How this Part answers RQ1},
  coltitle=gray!60!black,
  attach boxed title to top left={yshift=-2mm, xshift=6mm},
  boxed title style={colback=gray!15, arc=2pt, boxrule=0.6pt, colframe=gray!40},
]
RQ1 is a \emph{design} question, so it is answered by an artefact: a governance framework plus an adoption model. The Part runs \textbf{Foundations} (define ``safe adoption'' and anchor it to standards) $\rightarrow$ the \textbf{four pillars} taken directly from the question (governance, human-in-the-loop, audit logging, risk management) $\rightarrow$ \textbf{adoption, measurement and business case} $\rightarrow$ a \textbf{synthesis} that traces the framework back to the success criteria.\\[4pt]
\textit{Migration note: the earlier shallow proposal sections (Proposed Solution $\ldots$ Business Case) will be deepened and merged into the matching subsections below, then removed.}
\end{tcolorbox}

% ============================================================
\section{Framing and Foundations}

\subsection{From Question to Artefact}
% FILL: research logic — RQ1 is design-oriented; the deliverable that answers it is a framework + adoption model. State Part II structure and map each pillar to a clause of the question.
\textit{[To be developed: research logic and structure of Part II.]}

\subsection{Defining ``Safe Adoption''}
% FILL: operationalize "safe" as success criteria: (1) no unauthorized disclosure, (2) regulatory compliance (GDPR 25/32), (3) accountability & auditability, (4) preserved human oversight (no unsupervised high-risk release), (5) sustainable adoption. Define failure conditions. This becomes the yardstick reused in Synthesis.
\textit{[To be developed: operational success criteria for safe adoption.]}

\subsection{Regulatory and Standards Baseline}
% FILL: GDPR Art.25 (by design/default) & Art.32 (security of processing); ICO reprimand expectations; NIST AI RMF (Govern/Map/Measure/Manage); ISO/IEC 27001 (ISMS) & ISO/IEC 42001 (AI management system); PMI/PMBOK knowledge areas. End with a mapping table: requirement <-> standard <-> where addressed.
\textit{[To be developed: standards baseline + mapping table.]}
% TODO table: Requirement | Standard / clause | Addressed in (pillar / section)

\subsection{Framework Architecture}
% FILL: one-page overview + diagram — a governance "spine" supporting the 4 pillars, feeding an adoption lifecycle; restate the guiding principle "AI suggests, humans approve, the organization remains accountable."
\textit{[To be developed: framework overview + architecture diagram.]}

% ============================================================
\section{Governance (Pillar I)}

\subsection{Operating Model and Oversight}
% FILL: governance owners (Document Owner, Legal, DPO/Privacy, IT Security, Manager) + oversight forum; who owns the framework end to end.
\textit{[To be developed.]}

\subsection{GenAI-Assisted Redaction Policy}
% FILL: purpose, scope, principles, prohibited uses, vendor & data-handling rules (links to the "data exposure to external AI vendor" risk).
\textit{[To be developed.]}

\subsection{Decision Rights and Approval Gates}
% FILL: who may authorize release at each risk tier; gates that block automatic publication of high-risk documents.
\textit{[To be developed.]}

\subsection{RACI Matrix}
% FILL: migrate + deepen the existing RACI (current "Roles and Responsibilities" section); justify each A/R assignment.
\textit{[To be developed: migrate from existing RACI section.]}

% ============================================================
\section{Human-in-the-Loop Controls (Pillar II)}

\subsection{Document Classification and Risk Scoring}
% FILL: data categories (personal, financial, legal, confidential); how the risk score is computed (sensitivity x exposure x audience).
\textit{[To be developed.]}

\subsection{Risk-Tier to Control Matrix}
% FILL: migrate + deepen the existing HITL table (risk level -> automation level -> required human control); keep "very high = no automatic release".
\textit{[To be developed: migrate from existing HITL table.]}

\subsection{Review and Escalation Workflow}
% FILL: refine the 11-step workflow into swimlanes with control gates; mandatory review and escalation for high / very-high risk.
\textit{[To be developed: deepen the 11-step workflow.]}

\subsection{Reviewer Guidance and Automation Bias}
% FILL: what the human checks; mitigating over-reliance on AI (links to the "overreliance on AI" risk); reviewer competence / training requirement.
\textit{[To be developed.]}

% ============================================================
\section{Audit Logging and Accountability (Pillar III)}

\subsection{Purpose and Principles}
% FILL: immutability, completeness, traceability; the log as evidence of "appropriate technical and organizational measures".
\textit{[To be developed.]}

\subsection{Audit Log Schema}
% FILL: migrate + deepen the existing required fields; add retention period and access controls.
\textit{[To be developed: migrate from existing audit log section.]}

\subsection{Demonstrable Compliance}
% FILL: map log fields to GDPR accountability (Art.5(2)) and ICO expectations; how the log would answer a regulator.
\textit{[To be developed.]}

\subsection{Audit Review Process}
% FILL: review cadence, who reviews (IT Security / DPO), exception handling.
\textit{[To be developed.]}

% ============================================================
\section{Risk Management (Pillar IV)}

\subsection{Risk Management Approach}
% FILL: PMBOK-aligned process — identify -> qualitative/quantitative analysis -> response -> monitor & control.
\textit{[To be developed.]}

\subsection{Risk Register}
% FILL: migrate + deepen the existing register; add risk categories and root causes.
\textit{[To be developed: migrate from existing risk register.]}

\subsection{Qualitative and Quantitative Analysis}
% FILL: probability x impact scoring; risk matrix / heat map; optional expected monetary value (EMV) tied to the case (ICO sanction, remediation, reputational cost).
\textit{[To be developed: add heat map + EMV.]}

\subsection{Risk Responses and Residual Risk}
% FILL: avoid / reduce / transfer / accept; map mitigations to Pillars I-III; state residual risk + monitoring.
\textit{[To be developed.]}

% ============================================================
\section{Adoption, Measurement and Business Case}

\subsection{Adoption Roadmap}
% FILL: phased rollout Assess -> Design -> Pilot -> Scale (PM lifecycle; note iterative/Agile pilot from the course). Milestones and gate criteria.
\textit{[To be developed.]}

\subsection{Change Management and Training}
% FILL: stakeholder engagement, communication plan, training before go-live (links to the "lack of staff training" risk).
\textit{[To be developed.]}

\subsection{KPIs and Control Effectiveness}
% FILL: migrate + deepen the existing KPIs; split leading vs lagging; add control-effectiveness metrics.
\textit{[To be developed: migrate from existing KPI section.]}

\subsection{Maturity Model}
% FILL: 5-level redaction-governance maturity (ad hoc -> optimized) to position adoption progress.
\textit{[To be developed.]}

\subsection{Business Case}
% FILL: migrate + deepen the existing cost scenario; reframe as risk-driven ROI / cost of inaction anchored to the case.
\textit{[To be developed: migrate from existing business case.]}

% ============================================================
\section{Synthesis: Answering RQ1}

\subsection{How the Framework Achieves Safe Adoption}
% FILL: trace the framework back to the 5 success criteria ("Defining Safe Adoption") and the standards mapping; give the explicit answer to RQ1.
\textit{[To be developed.]}

\subsection{Limitations and Assumptions}
% FILL: scope limits, assumptions, what is not guaranteed.
\textit{[To be developed.]}

\subsection{Bridge to RQ2}
% FILL: set up applying the framework to the Post Office Horizon case (RQ2).
\textit{[To be developed.]}

% ============================================================
% SECTION 16
% ============================================================
\section{Conclusion}

The Post Office Horizon case shows that document redaction failures are not only technical errors. They can also be failures of governance, accountability and process control.

A GenAI-assisted redaction framework could reduce this risk by supporting sensitive data detection, risk scoring, redaction suggestions and pre-publication checks. However, GenAI should not replace human responsibility.

\begin{principlebox}[Final Takeaway]
The safest approach is a controlled workflow where AI supports the process, humans approve high-risk decisions and every action is recorded through an audit log.
\end{principlebox}

% ============================================================
% REFERENCES
% ============================================================
\section{Preliminary References}

\begin{thebibliography}{9}

\bibitem{ico2025}
Information Commissioner's Office.
\textit{Post Office reprimanded over Horizon IT scandal victims' entirely preventable data breach}.
2025.
Available at: \url{https://ico.org.uk/about-the-ico/media-centre/news-and-blogs/2025/12/post-office-reprimanded-over-horizon-it-scandal-victims-entirely-preventable-data-breach/}

\bibitem{icoreprimand2025}
Information Commissioner's Office.
\textit{Post Office Limited Reprimand}.
2025.
Available at: \url{https://ico.org.uk/media2/sroaikut/post-office-limited-reprimand-20251202.pdf}

\bibitem{guardian2024}
The Guardian.
\textit{Post Office accidentally leaks names and addresses of wrongfully convicted operators}.
2024.

\bibitem{gdpr}
European Union.
\textit{General Data Protection Regulation, Articles 25 and 32}.

\end{thebibliography}

\end{document}
